Kết quả từ mô hình tối ưu hóa Copula-CVaR cho thấy một sự dịch chuyển quan trọng trong việc tiếp cận quản trị rủi ro: sự an toàn của danh mục không đến từ việc loại bỏ hoàn toàn biến động, mà đến từ việc kiểm soát những tổn thất cực đoan nhất. Bằng cách ghi nhận hiện tượng đuôi dày và sự phụ thuộc phi tuyến, mô hình đã thành công trong việc nhận diện những khiếm khuyết của phương pháp Markowitz truyền thống. Việc ngưỡng rủi ro $VaR$ được ép xuống mức thấp hơn chứng minh rằng, khi thuật toán nhận thức được xác suất xảy ra các đợt bán tháo đồng loạt, nó sẽ tự động cơ cấu lại tỷ trọng để tạo ra một danh mục có sức đề kháng mạnh mẽ hơn trước các cú sốc hệ thống.